\begin{center}\begin{tikzpicture}[font=\small,>=stealth,line join=round]
\draw[gray!40,rounded corners=3pt] (0,0) rectangle (14.2,8.8);\node at (7.1,8.200000000000001) {证明三：把中心移到原点，转成奇函数关系};
\node at (3.1500000000000004,7.3) {$f(x)=3+\sin\pi(x-\frac12)$};\begin{scope}[shift={(1.6800000000000002,3.4499999999999997)},x=0.9800000000000001cm,y=0.5937499999999999cm]\draw[->,gray] (-1,0)--(4,0) node[right]{$x$};\draw[->,gray] (0,-1.6)--(0,4.8) node[above]{$y$};\draw[teal,thick,solid,domain=-0.7:3.7,samples=180] plot(\x,{3+sin(deg(pi*(\x-.5)))});\fill[orange!85!black] (0.5,3) circle(2pt);\fill[orange!85!black] (1.5,3) circle(2pt);\draw[gray,dashed] (-.8,3)--(3.8,3);\end{scope}\node[align=center,text width=5.8cm] at (3.1500000000000004,1.7) {中心 $(\frac12,3)$\\原来的周期为2};\node at (9.95,7.3) {$h(t)=f(\frac12+t)-3$};\begin{scope}[shift={(8.48,3.4499999999999997)},x=0.9800000000000001cm,y=0.5937499999999999cm]\draw[->,gray] (-1,0)--(4,0) node[right]{$t$};\draw[->,gray] (0,-1.6)--(0,4.8) node[above]{$y$};\draw[orange!85!black,thick,solid,domain=-0.7:3.7,samples=180] plot(\x,{sin(deg(pi*\x))});\fill[orange!85!black] (0,0) circle(2pt);\fill[orange!85!black] (1,0) circle(2pt);\end{scope}\node[align=center,text width=5.8cm] at (9.95,1.7) {中心变为 $(0,0)$\\$h=\sin\pi t$，周期仍为2};
\node[text width=12.6cm,align=center] at (7.1,.65) {两图同尺度；坐标替换和减去中心高度，不改变横向重复的距离。};
\end{tikzpicture}\end{center}
